% ===================================================================
%  بخش ۵ — نتیجه‌گیری  (بازنویسی بر پایه‌ی اجرای دوم)
% ===================================================================
\section{نتیجه‌گیری}
\label{sec:conclusion}

در این پژوهش، پاسخ یک گسیل‌گر دوقطبی در فاصله‌ی پنج نانومتری از نانومیله‌ی
طلای کلاهک‌دار با روش 3D-FDTD در ده اجرای کنترل‌شده و اسکریپت‌محور
بررسی شد (جدول~\ref{tab:cases}: A, B, C, D, D$_\perp$, M3, M1.5, G3, G10, G20). برای دوقطبی طولی، فاکتور پورسل در $\nm{610}$ به $1.5\times10^{3}$
و تقویت توان تابشی در $\nm{616}$ به $1.2\times10^{2}$ می‌رسد. طیف کامل
بازده نشان می‌دهد که در قله‌ی پورسل تنها $7.6\%$ و در قله‌ی تابشی
$8.6\%$ توان تابیده می‌شود و بیش از $92\%$ به کانال غیرتابشی می‌رود؛
گزارش فاکتور پورسل به‌تنهایی برای ارزیابی نانوآنتن کافی نیست.

چهار کنترل نقش جهت و شکل را از هم جدا کردند. چرخاندن دوقطبی به راستای
عرضی روی همان نوک، فاکتور پورسل در تشدید را $41$ برابر و توان تابشی را
$360$ برابر کاهش داد. اما نتیجه‌ی تعیین‌کننده از مقایسه با کره‌ی هم‌حجم
با دو جهت‌گیری شعاعی و مماسی به دست آمد: نسبت جهتی کره در کل طیف
تقریباً ثابت و نزدیک $2.4$ است، و نسبت جهتی نانومیله در خارج از تشدید
دقیقاً همین مقدار است ولی در تشدید طولی یک تا دو مرتبه‌ی بزرگی می‌جهد.
بنابراین ناهمسانگردی تانسور گرین نانومیله دو سهم دارد: سهم هندسی
نزدیک‌سطحی از مرتبه‌ی $2$ که هر جسم فلزی دارد، و سهم تشدیدی که تنها به
مد طولی محور ممتاز تعلق دارد و توصیف اسکالر یا میانگین‌گیری جهتی آن را
حذف می‌کند. پویش گاف نیز نشان داد که با افزایش گاف از $5$ به $\nm{20}$
مکان تشدید ثابت می‌ماند، قله‌ی تابشی از $119$ به $10$ می‌افتد و بیشینه‌ی
بازده طیفی از $19\%$ به $74\%$ می‌رسد.

اعتبار عددی با سه آزمون سنجیده شد. کنترل فضای آزاد صحت شار را تأیید کرد
ولی حساسیت $\pm13\%$ نرمال‌سازی را نشان داد؛ مقایسه با حل می برای کره،
مکان قله‌ها را با اختلاف $\nm{2}$ بازتولید و کم‌برآوردی $8$ تا $14\%$
مش $\nm{2}$ را کمّی کرد؛ و آزمون همگرایی مش ($3$، $2$ و $\nm{1.5}$)
آشکار ساخت که مکان قله‌ی تابشی ($616\pm\nm{4}$) و بازده در آن پایدارند،
ولی ارتفاع قله‌ها به دلیل قرارگیری زیرسلولی چشمه در گاف پنج‌نانومتری
تا $\pm40\%$ تغییر می‌کند. به همین دلیل، ادعاهای این مقاله آگاهانه بر
مکان تشدیدها، بازده و نسبت‌های میان حالت‌های هم‌مش استوار شده‌اند و
ارتفاع مطلق قله‌ها با این عدم‌قطعیت گزارش می‌شود.

دامنه‌ی ادعای مقاله محدود به همین سامانه‌ی کلاسیک و موضعی است. نتایج به
معنی نامعتبر بودن ضرایب اینشتین برای اتم آزاد یا نظریه‌ی می برای کره
نیستند؛ بلکه نشان می‌دهند در محیط ناهمسانگرد و اتلافی، نرخ کل، نرخ تابشی
و جهت‌گیری دوقطبی باید جداگانه مشخص شوند. گام بعدی روشن و مشخص است:
اجرای حالت اصلی با مش $\le\nm{1}$ و چشمه‌ی دقیقاً روی گره (یا مش
ناهمسانگرد ریز در راستای گاف) تا ارتفاع قله‌ها نیز همگرا شود؛ اجرای
مرجع فضای آزاد با همان مش و مختصات؛ و گزارش برازش داده‌ی طلا. توسعه‌ی
پایه‌ی مدی متناسب با تقارن استوانه‌ای
\cite[pp.~123--131]{maier2007plasmonics} و بررسی پاسخ غیرموضعی موضوع
کار آینده‌اند.

از دید کاربردی، تقویت بیش از صدبرابری توان تابشی نشان می‌دهد نانومیله‌ی
طلا برای کاربردهایی که نرخ گسیل و شدت سیگنال مهم‌تر از بازده‌اند —
حسگری فلورسانسی و طیف‌سنجی تقویت‌شده — مفید است؛ اما برای چشمه‌ی
تک‌فوتونی، که به بازده جمع‌آوری بالا نیاز دارد، پویش گاف حاضر
(شکل~\ref{fig:gap}) نشان می‌دهد گاف $10$ تا $\nm{20}$ با تقویت تابشی
$10$ تا $50$ برابری و بازده رو به بهبود، مصالحه‌ی معقول‌تری است
\cite{mohammadi2008gold, koenderink2017single}. تحقق تجربی چنین
ساختاری نیز به کنترل اندازه و شکل نانومیله در فرآیند سنتز وابسته است
\cite{nikoobakht2003preparation}.
